\part*{優先度B：知識がないと時間がかかる}
\addcontentsline{toc}{part}{優先度B：知識がないと時間がかかる}

\chapter{Kaplan--Meier 推定}\label{ch:m04}

\begin{priobox}{B}
「カプラン・マイヤー法によって描け」（\kakomon{2016}{3}〔1〕，\kakomon{2021}{1}〔2〕）は推定式が与えられないが，
準1級の範囲で式自体は易しい．落としやすいのは\textbf{同時点の打ち切りの扱い}，\textbf{リスク集合の数え間違い}，\textbf{打ち切りで曲線を下げる}誤りである．
\kakomon{2019}{1}〔2〕では与えられた Nelson--Aalen 型推定値が KM 以上であることの証明が問われた．本編の第9章9.4・9.6節．
\end{priobox}

% ============================================================
\section{概要}\label{sec:m04-gaiyo}
相異なるイベント時点を $t_{(1)}<\dots<t_{(r)}$，時点 $t_{(j)}$ の直前に観察中の人数（リスク集合の大きさ）を $n_j$，$t_{(j)}$ のイベント数を $d_j$ とする．
\begin{pointbox}[覚える式]
\[
\hat S(t)=\prod_{j:\,t_{(j)}\le t}\Bigl(1-\frac{d_j}{n_j}\Bigr)\quad\text{（Kaplan--Meier 推定量）}
\]
\begin{itemize}
\item 打ち切り時点では $\hat S$ は変わらず，次のリスク集合が減るだけ．同時点ではイベントが先（その時点の打ち切り例は $n_j$ に含める）．
\item 中央生存時間 $=\hat S(t)\le0.5$ となる最小の $t$．
\item Greenwood（近似）：$\widehat\V[\hat S(t)]\approx\hat S(t)^2\sum_{t_{(j)}\le t}\dfrac{d_j}{n_j(n_j-d_j)}$．
\item Nelson--Aalen：$\tilde H(t)=\sum_{t_{(j)}\le t}d_j/n_j$，$\tilde S(t)=e^{-\tilde H(t)}\ge\hat S(t)$（$e^{-x}\ge1-x$ による）．
\end{itemize}
\end{pointbox}

% ============================================================
\section{詳細}\label{sec:m04-shosai}
\subsection{式の意味}
$t_{(j)}\le t<t_{(j+1)}$ なら，イベントのない区間では生存確率は減らないので
\[
S(t)=\prod_{k\le j}\Prob\bigl(T>t_{(k)}\mid T\ge t_{(k)}\bigr)
\]
と条件付き確率の積に分解できる．各因子を「時点 $t_{(k)}$ のリスク集合 $n_k$ 人のうち生き残った割合」$1-d_k/n_k$ で推定したものが KM 推定量である．
打ち切り例は打ち切り時点まではリスク集合に入り，その後の分母から外れる．打ち切り例を最初から除いたり，打ち切りをイベントとみなしたりすると偏る．
前提は独立な打ち切り（打ち切られた人のその後のハザードが，同じ時点でリスク集合に残っている人と同じ）である．

\subsection{表の作り方と曲線}
\begin{enumerate}
\item 観察時間を昇順に並べる．打ち切り指示変数の向き（1が打ち切りかイベントか）を確認する．
\item 各イベント時点で $n_j$（その時点以上の観察時間をもつ人数）と $d_j$ を数える．
\item 前の $\hat S$ に $(1-d_j/n_j)$ を掛ける．
\item 曲線は右連続な階段関数で，イベント時点でのみ下がる．打ち切り時点には $+$ 印を付ける．
      最長の観察が打ち切りなら曲線は0に到達しない．最長の観察がイベントなら $n=d$ で0になる．
\end{enumerate}
過去問の解答の表には「死亡率」の列（$1-\hat S$）や打ち切り時点の行も含まれることがある．$1-\hat S(t)$ は累積イベント割合（累積発生率）の推定値で，\kakomon{2021}{1}〔2〕はこれを描かせた．

\subsection{Greenwood の公式}
$\hat p_j=1-d_j/n_j$ とすると $\log\hat S=\sum_j\log\hat p_j$．$n_j$ を与えると $n_j-d_j\sim\Bin(n_j,p_j)$ で，
(i) 異なる時点の $\hat p_j$ を無相関とみなし，(ii) デルタ法で $\V[\log\hat p_j]\approx(1-p_j)/(n_jp_j)$ として $p_j$ に $\hat p_j$ を代入すると $d_j/\{n_j(n_j-d_j)\}$．
最後に $e^x$ のデルタ法で $\V[\hat S]\approx\hat S^2\V[\log\hat S]$．厳密な分散ではなく漸近近似である．
$n_j=d_j$（最後のイベントでリスク集合が尽きる）と定義できない．
\kakomon{2016}{3} の解答の「標準偏差」の列はこの公式による（公式自体は問われていない）．

\subsection{Nelson--Aalen 推定量との比較}
$d_j/n_j$ は時点 $t_{(j)}$ の離散ハザード（その時点でのイベントの割合）の推定値である．
\begin{itemize}
\item KM は離散時間の関係 $S=\prod(1-\text{離散ハザード})$ で生存関数を作る．
\item Nelson--Aalen は累積ハザードを $\tilde H=\sum d_j/n_j$ と推定し，連続時間の関係 $S=e^{-H}$ で $\tilde S=e^{-\tilde H}=\prod e^{-d_j/n_j}$ とする．
\end{itemize}
\begin{meidai}[$\tilde S\ge\hat S$（\kakomon{2019}{1}〔2〕）]\label[meidai]{prop:m04-na}
任意の $t$ で $\tilde S(t)\ge\hat S(t)$．
\end{meidai}
\begin{proof}
$g(x)=e^{-x}-(1-x)$ は $g(0)=0$，$g'(x)=1-e^{-x}$ が $x<0$ で負，$x>0$ で正なので $x=0$ で最小値0をとり，すべての $x$ で $e^{-x}\ge1-x$．
$0\le d_j/n_j\le1$ なので各因子で $e^{-d_j/n_j}\ge1-d_j/n_j\ge0$．非負の数どうしの積なので大小関係が保たれ，$\tilde S(t)\ge\hat S(t)$．
\end{proof}
リスク集合が大きく $d_j/n_j$ が小さい時点では両者はほぼ等しく，右端（$n_j$ が小さい）で差が開く．最後のイベントで $d_j=n_j$ なら KM は0，$\tilde S$ は0にならない．

\begin{warnbox}
\begin{itemize}
\item 同じ時点の打ち切りをイベントより先に処理し，リスク集合を1減らす．
\item 打ち切り時点で曲線を下げる．
\item 中央値を「$\hat S=0.5$ を線形補間した時点」とする（定義は $\hat S(t)\le0.5$ となる最小の $t$）．
\item 打ち切り指示変数の向きを読み違える（2016・2019年は1が打ち切り，2018・2021年は1がイベント）．
\end{itemize}
\end{warnbox}

% ============================================================
\section{出題例}\label{sec:m04-shutsudai}
\begin{shutsudai}{\kakomon{2016}{3}〔1〕：2群の KM 曲線}
群1は最初の3人が打ち切り（3，5，7週）で，曲線は9週まで1のまま，9週で $1/2$，10週で0．
群2は1週 $4/5=0.8$，2週 $0.6$，4週の打ち切りを挟んで6週 $0.6\times1/2=0.3$，8週で0．
「1：打ち切り，0：イベント」の向きに注意．
\end{shutsudai}
\begin{shutsudai}{\kakomon{2019}{1}〔1〕〔2〕：KM 表と Nelson--Aalen}
5例（10，13，$18^+$，19，$23^+$ 週）で $\hat S(10)=0.8$，$\hat S(13)=0.6$，$\hat S(19)=0.3$（19週のリスク集合は2人）．
〔2〕は $\tilde S(t)=\prod\exp(-d_k/n_k)$ が与えられ，$\tilde S\ge\hat S$ を示す（\cref{prop:m04-na}）．
〔3〕以降は境界内平均生存時間（RMST）．KM による RMST は曲線下の長方形の和 $10\times1+3\times0.8+6\times0.6+1\times0.3=16.3$（$\tau=20$）．
\end{shutsudai}
\begin{shutsudai}{\kakomon{2021}{1}〔2〕：$1-\hat S(t)$ のグラフ}
8例（打ち切り4，6，8，16年，イベント10，12，14，18年）．$\hat S$ は10年0.8，12年0.6，14年0.4，18年0．縦軸 $1-\hat S$ で描く．
この後，打ち切りの一部が競合イベントだった場合の比較に続く（第5章）．
\end{shutsudai}

% ============================================================
\section{例題}\label{sec:m04-reidai}
\begin{reidai}[\kakomon{2019}{1}〔1〕〔2〕・\kakomon{2021}{1}〔2〕型：KM 表・曲線・Nelson--Aalen]\label{reidai:m04-1}
6人の生存時間データ（週）を示す（イベントの有無は 0：打ち切り，1：イベント）．
\begin{center}\small
\begin{tabular}{@{}ccccccc@{}}
\toprule
被験者 & 1 & 2 & 3 & 4 & 5 & 6\\
\midrule
観察時間（週） & 11 & 22 & 8 & 11 & 17 & 11\\
イベント & 1 & 0 & 1 & 0 & 1 & 1\\
\bottomrule
\end{tabular}
\end{center}
\begin{enumerate}[label=〔\arabic*〕]
\item KM 推定値の表（$t_{(j)}$，$n_j$，$d_j$，$\hat S$）を作り，生存曲線を描け．中央生存時間を求めよ．
\item Greenwood の公式で $\hat S(11)$ の標準誤差を求めよ．
\item 各イベント時点の Nelson--Aalen 型推定値 $\tilde S(t)=\prod_{t_{(k)}\le t}\exp(-d_k/n_k)$ を求めよ．また任意の $t$ で $\tilde S(t)\ge\hat S(t)$ となることを示せ．
\item 累積イベント割合 $1-\hat S(t)$ の推定値を各イベント時点で求めよ．
\end{enumerate}
\end{reidai}

% ============================================================
\section{例題の詳しい解説}\label{sec:m04-kaisetsu}
\begin{kaitou}{reidai:m04-1}
〔1〕「1：イベント」の向きで読むと，昇順に $8,\ 11,\ 11,\ 11^+,\ 17,\ 22^+$．
11週にはイベント2件と打ち切り1件が同時にある．\textbf{イベントが先}なので，11週のリスク集合は打ち切りの人も含めて5人．
\begin{center}\small
\begin{tabular}{@{}cccccc@{}}
\toprule
$t_{(j)}$ & $n_j$ & $d_j$ & $1-d_j/n_j$ & $\hat S(t)$ & $1-\hat S(t)$\\
\midrule
8 & 6 & 1 & $5/6$ & $0.8333$ & $0.1667$\\
11 & 5 & 2 & $3/5$ & $0.8333\times0.6=0.5000$ & $0.5000$\\
17 & 2 & 1 & $1/2$ & $0.5000\times0.5=0.2500$ & $0.7500$\\
\bottomrule
\end{tabular}
\end{center}
（17週のリスク集合は，11週のイベント2件と打ち切り1件で5人から3人減って2人．）
\begin{center}
\begin{tikzpicture}
\begin{axis}[width=90mm,height=55mm,xmin=0,xmax=24,ymin=0,ymax=1.05,
  xlabel={時間（週）},ylabel={$\hat S(t)$},xtick={0,8,11,17,22},ytick={0,0.25,0.5,0.8333,1},
  yticklabels={0,0.25,0.5,0.833,1},grid=major,grid style={gray!20},tick label style={font=\scriptsize}]
\addplot[very thick,bkblue,const plot] coordinates {(0,1)(8,0.8333)(11,0.5)(17,0.25)(24,0.25)};
\addplot[only marks,mark=+,mark size=3pt,thick,bkorange] coordinates {(11,0.5)(22,0.25)};
\addplot[dashed,gray] coordinates {(0,0.5)(24,0.5)};
\end{axis}
\end{tikzpicture}
\end{center}
最長の観察（22週）が打ち切りなので曲線は0.25で止まる．
中央生存時間は $\hat S(t)\le0.5$ となる最小の $t$ で，$\hat S(11)=0.5$ なので\textbf{11週}（ちょうど0.5でも「以下」に入る）．

〔2〕$\sum_{t_{(j)}\le11}\dfrac{d_j}{n_j(n_j-d_j)}=\dfrac1{6\times5}+\dfrac2{5\times3}=0.03333+0.13333=0.16667$．
$\widehat\V[\hat S(11)]\approx0.5^2\times0.16667=0.04167$，$\SE=\sqrt{0.04167}=0.204$．6人の小標本なので，この近似の SE は目安にすぎない．

〔3〕$\tilde H$ は $d_j/n_j$ の累積和で，$e^{-\tilde H}$ は付表から引く：
\begin{center}\small
\begin{tabular}{@{}ccccc@{}}
\toprule
$t_{(j)}$ & $d_j/n_j$ & $\tilde H(t)$ & $\tilde S(t)=e^{-\tilde H}$ & $\hat S(t)$\\
\midrule
8 & $1/6=0.1667$ & $0.1667$ & $0.8465$ & $0.8333$\\
11 & $2/5=0.4000$ & $0.5667$ & $0.5674$ & $0.5000$\\
17 & $1/2=0.5000$ & $1.0667$ & $0.3442$ & $0.2500$\\
\bottomrule
\end{tabular}
\end{center}
すべての時点で $\tilde S\ge\hat S$ で，差は右へ行くほど大きい．
証明：$g(x)=e^{-x}-(1-x)$ とおくと $g'(x)=1-e^{-x}$ は $x<0$ で負，$x>0$ で正なので，$g$ は $x=0$ で最小値 $g(0)=0$ をとり，$e^{-x}\ge1-x$．
$x=d_k/n_k\in[0,1]$ を代入すると $\exp(-d_k/n_k)\ge1-d_k/n_k\ge0$．
$t_{(j)}\le t<t_{(j+1)}$ では $\tilde S(t)$ と $\hat S(t)$ はそれぞれ $k=1,\dots,j$ の因子の積で，各因子が非負かつ大小関係が揃っているので積でも $\tilde S(t)\ge\hat S(t)$．
$t<t_{(1)}$ では両方とも1．

〔4〕$1-\hat S$ は8週 $0.167$，11週 $0.500$，17週 $0.750$（〔1〕の表の右端の列）．17週以降は0.75のまま．
\end{kaitou}

\begin{summarybox}
\begin{itemize}
\item $\hat S(t)=\prod_{t_{(j)}\le t}(1-d_j/n_j)$．同時点ではイベントが先．打ち切りで曲線は下げない．
\item 中央値は $\hat S\le0.5$ の最小時点．最長観察が打ち切りなら曲線は0にならない．
\item Greenwood は近似：$\hat S^2\sum d_j/\{n_j(n_j-d_j)\}$．Nelson--Aalen $\tilde S=e^{-\sum d_j/n_j}\ge\hat S$．
\end{itemize}
\end{summarybox}
